% ============================================================
%  commands.tex — Custom macros for ECON 510 (Econometrics)
%  Adapted from adv-micro-psu, expanded for asymptotic theory
% ============================================================

\usepackage{amsmath}

% Environment shortcuts
\newcommand{\ba}{\begin{aligned}}
\newcommand{\ea}{\end{aligned}}
\newcommand{\bc}{\begin{cases}}
\newcommand{\ec}{\end{cases}}

% Number sets
\newcommand{\RR}{\mathbb{R}}
\newcommand{\QQ}{\mathbb{Q}}
\newcommand{\ZZ}{\mathbb{Z}}
\newcommand{\NN}{\mathbb{N}}

% Calligraphic
\renewcommand{\P}{\mathcal{P}}
\renewcommand{\S}{\mathcal{S}}
\newcommand{\T}{\mathcal{T}}
\newcommand{\A}{\mathcal{A}}
\newcommand{\B}{\mathcal{B}}
\newcommand{\C}{\mathcal{C}}
\renewcommand{\L}{\mathcal{L}}
\newcommand{\F}{\mathcal{F}}
\newcommand{\W}{\mathcal{W}}

% Bold (use distinct names to avoid clash with calligraphic)
\newcommand{\vb}[1]{\mathbf{#1}}
\newcommand{\bX}{\mathbf{X}}
\newcommand{\bY}{\mathbf{Y}}
\newcommand{\bW}{\mathbf{W}}
\newcommand{\bZ}{\mathbf{Z}}
\newcommand{\bU}{\mathbf{U}}

% Greek
\newcommand{\ve}{\varepsilon}
\newcommand{\s}{\sigma}
\renewcommand{\l}{\lambda}
\newcommand{\g}{\gamma}
\newcommand{\Th}{\Theta}
\renewcommand{\th}{\theta}    % overwrite built-in thorn (\th)
\newcommand{\Om}{\Omega}
\newcommand{\om}{\omega}
\newcommand{\Gm}{\Gamma}

% Delimiters
\newcommand{\br}[1]{\left[#1\right]}
\newcommand{\pr}[1]{\left(#1\right)}
\renewcommand{\brace}[1]{\left\{#1\right\}}
\newcommand{\abs}[1]{\left|#1\right|}
\newcommand{\norm}[2][]{\left\| #2 \right\|_{#1}}
\newcommand{\floor}[1]{\left\lfloor #1 \right\rfloor}
\newcommand{\ceil}[1]{\left\lceil #1 \right\rceil}

% Operators
\DeclareMathOperator*{\argmax}{arg\,max}
\DeclareMathOperator*{\argmin}{arg\,min}
\DeclareMathOperator{\plim}{plim}
\DeclareMathOperator{\rk}{rk}
\DeclareMathOperator{\tr}{tr}
\DeclareMathOperator{\diag}{Diag}
\DeclareMathOperator{\sgn}{sgn}
\DeclareMathOperator{\Corr}{Corr}

% Probability/Expectation
% Usage: \E{x \given y}, \Pr{A}, \Var{x}, \Cov{x}{y}.
% For starred bootstrap versions, write \mathbb{E}^*, \mathrm{P}^*, etc. directly.
\newcommand{\given}{\,|\,}   % plain bar: works in subscripts and \pr{} alike
\renewcommand{\Pr}[1]{\operatorname{P}\!\left(#1\right)}
\newcommand{\E}[2][]{\mathbb{E}_{#1}\!\pr{#2}}
% \Var: parentheses (econometric convention). Accepts both \Var{X} and \Var(X).
% Calling \Var{X} -> Var(X); calling \Var(X) (explicit parens) also works since
% the macro renders Var with no implicit delimiter when called as a no-arg op.
\newcommand{\Var}[2][]{\operatorname{Var}_{#1}\!\pr{#2}}
\newcommand{\Cov}[2]{\operatorname{Cov}\!\pr{#1,#2}}

% Convergence
\newcommand{\dto}{\overset{d}{\to}}
\newcommand{\pto}{\overset{p}{\to}}
\newcommand{\asto}{\overset{a.s.}{\to}}
\newcommand{\dteq}{\overset{d}{=}}
\newcommand{\iid}{\overset{i.i.d.}{\sim}}

% Asymptotic order (the user's weak point — treat clearly)
\newcommand{\Op}{O_{p}}      % O_p
\newcommand{\op}{o_{p}}      % o_p
\newcommand{\rtn}{\sqrt{n}}  % \sqrt{n} shortcut

% Calculus
\renewcommand{\d}{\text{d}}
\newcommand{\dd}[2]{\frac{\d #1}{\d #2}}
\newcommand{\pd}[2]{\frac{\partial #1}{\partial #2}}
\newcommand{\inv}[1]{{#1}^{-1}}

% Logic
\newcommand{\se}{\subseteq}
\newcommand{\bs}{\setminus}
\newcommand{\st}{\text{ s.t. }}
\newcommand{\mif}{\text{ if }}
\newcommand{\miff}{\;\text{iff}\;}

% Hat / bar / tilde shortcuts (used everywhere in econometrics)
\newcommand{\hb}{\widehat{\beta}}
\newcommand{\hth}{\widehat{\theta}}
\newcommand{\tth}{\widetilde{\theta}}

% Sums / products
\newcommand{\sumi}{\sum_{i=1}^n}
\newcommand{\prodi}{\prod_{i=1}^n}

% Highlighting
\newcommand{\match}[1]{\textcolor{ForestGreen}{$\boldsymbol{#1}$}}
\newcommand{\hl}[1]{\textcolor{red!70!black}{\textbf{#1}}}

% ----------------------------------------------------------------
% Proof tier tags (place at the start of a proof body)
% ----------------------------------------------------------------
%   \reproduce  — must memorize and reproduce on the exam
%   \structure  — know the steps; details may be fudged
%   \intuition  — included for understanding only; do not memorize
%   \proofskip  — included for completeness; safe to skip entirely
% (\skip is reserved by LaTeX, so we use \proofskip.)
\newcommand{\reproduce}{\textcolor{ForestGreen!70!black}{\textbf{[REPRODUCE — memorize this proof]}}\par\smallskip}
\newcommand{\structure}{\textcolor{orange!85!black}{\textbf{[STRUCTURE — know the steps, fudge details]}}\par\smallskip}
\newcommand{\intuition}{\textcolor{blue!70!black}{\textbf{[INTUITION ONLY — read once, do not memorize]}}\par\smallskip}
\newcommand{\proofskip}{\textcolor{gray!50!black}{\textbf{[SKIP — included for completeness, safe to skip]}}\par\smallskip}
